\documentclass[a4paper,twoside]{article}

\usepackage{morewrites}

\usepackage[svgnames,dvipsnames,x11names]{xcolor}
\usepackage{graphicx}
\usepackage[english]{babel}
\usepackage[colorlinks=true, linkcolor=NavyBlue, urlcolor=NavyBlue, citecolor=NavyBlue]{hyperref}
\usepackage[a4paper]{geometry}
\usepackage{ts-fonts}
\usepackage{amsmath}
\usepackage{float}
\usepackage{seqsplit}
\usepackage{ts-code}

\usepackage{lastpage}
\usepackage{refcount}
\makeatletter
\def\ps@plain{
  \let\@oddhead\@empty
  \let\@evenhead\@empty
  \def\@oddfoot{
    \hfil
    \ifnum\getpagerefnumber{LastPage}>1\relax
      \thepage
    \fi
    \hfil}
  \let\@evenfoot\@oddfoot
}
\makeatother

\title{Plugin API}
\author{}\date{}

\begin{document}

\maketitle

\thispagestyle{plain}
Plugins bundle opinionated HTML post-processors and asset helpers so you can
extend TeXSmith without patching the core read \(\rightarrow\) IR \(\rightarrow\) write pipeline.

\texttt{texsmith.plugins} exposes a namespace package populated by the MkDocs hook in
\texttt{docs/hooks/mkdocs\_hooks.py}, re-exporting the maintained plugin modules under
\texttt{texsmith.adapters.plugins}.

\section{Loading plugins}\label{loading-plugins}

\begin{code}{python}{}{}
\begin{Verbatim}[commandchars=\\\{\},breaklines, breakanywhere, commandchars=\\\{\}]
\PY{k+kn}{import}\PY{+w}{ }\PY{n+nn}{texsmith}\PY{n+nn}{.}\PY{n+nn}{plugins}\PY{n+nn}{.}\PY{n+nn}{snippet}  \PY{c+c1}{\PYZsh{} noqa: F401}

\PY{k+kn}{from}\PY{+w}{ }\PY{n+nn}{texsmith}\PY{+w}{ }\PY{k+kn}{import} \PY{n}{Document}\PY{p}{,} \PY{n}{convert\PYZus{}documents}

\PY{n}{bundle} \PY{o}{=} \PY{n}{convert\PYZus{}documents}\PY{p}{(}\PY{p}{[}\PY{n}{Document}\PY{o}{.}\PY{n}{from\PYZus{}markdown}\PY{p}{(}\PY{n}{Path}\PY{p}{(}\PY{l+s+s2}{\PYZdq{}}\PY{l+s+s2}{intro.md}\PY{l+s+s2}{\PYZdq{}}\PY{p}{)}\PY{p}{)}\PY{p}{]}\PY{p}{)}
\end{Verbatim}
\end{code}
\section{Authoring your own plugin}\label{authoring-your-own-plugin}

\begin{enumerate}
\item Create a module (e.g., \texttt{texsmith.plugins.acme}) exposing a Markdown
   extension and/or the HTML rewriting helpers your documents need.
\item Declare entry points or instruct consumers to \texttt{import texsmith.plugins.acme}
   before rendering.
\item Optionally provide a MkDocs plugin/hook so documentation builds load your
   plugin automatically.

\end{enumerate}
Keep plugin modules small and focused.

\section{Reference}\label{reference}

::: texsmith.plugins

::: texsmith.plugins.snippet

\end{document}
